\chapter{Conclusão}
\label{chp:conclusao}

\section{Considerações Finais}
    Este trabalho apresentou o desenvolvimento do VAPOREON, uma ferramenta de apoio ao pentest para iniciantes,
    com o objetivo de facilitar o aprendizado e a prática de testes de penetração. Através da integração de diversas 
    ferramentas de segurança em um ambiente virtualizado, o VAPOREON oferece uma interface gráfica intuitiva, 
    fluxos guiados e suporte por meio de LLM, contribuindo para a formação acadêmica e profissional de estudantes da área de
    TI. A avaliação realizada em ambiente controlado demonstrou a eficácia da ferramenta em proporcionar uma experiência de
    aprendizado prática

\section{Limitações do Estudo}
    Apesar dos resultados positivos, é importante reconhecer as limitações deste estudo. A avaliação foi conduzida 
    em um ambiente controlado, o que pode não refletir completamente as condições do mundo real. Além disso, a 
    ferramenta foi testada principalmente por usuários iniciantes, e a experiência de profissionais mais experientes 
    pode variar. A integração de LLMs, embora promissora, ainda apresenta desafios em termos de precisão e contextualização
    das respostas fornecidas.

\section{Trabalhos Futuros}
    Para aprimorar o VAPOREON e expandir suas funcionalidades, sugerem-se as seguintes direções para trabalhos futuros:
    \begin{itemize}
        \item Implementação de suporte a mais ferramentas de testes de penetração, ampliando a variedade de técnicas disponíveis para os usuários.
        \item Desenvolvimento de módulos específicos para testes de segurança em aplicações móveis e IoT, atendendo a demandas crescentes nessas áreas.
        \item Integração com plataformas de aprendizado online, permitindo que os usuários acessem conteúdos educacionais complementares diretamente pela interface do VAPOREON.
        \item Realização de estudos de usabilidade com um público mais amplo, incluindo profissionais experientes, para identificar oportunidades de melhoria na interface e nos fluxos guiados.
        \item Exploração de técnicas avançadas de inteligência artificial para oferecer suporte mais personalizado e contextualizado durante a utilização das ferramentas.
    \end{itemize}